\subsection{Actual projective model ideal functions and homogeneous coordinates}
\label{subsec:native-projective-model-ideals}

\begin{theorem}[Actual polynomial projective structure morphism is proper]
\label{projective:native-polynomial-structure-proper}
\lean{EnvelopeProperty.ProjBerk.homogeneousCoefficientZero_bijective,
EnvelopeProperty.ProjBerk.homogeneousCoefficientZeroIso,
EnvelopeProperty.ProjBerk.nativeStructureMap_isProper}
\leanok
\uses{projective:native-structure-morphism}
For every commutative coefficient ring $R$, the genuine constant
homomorphism from $R$ to the degree-zero homogeneous component of
$R[X_0,\ldots,X_n]$ is a whole ring isomorphism. The actual polynomial
projective structure morphism to $\operatorname{Spec}R$ is proper.
\end{theorem}
\begin{proof}
The original constant polynomial homomorphism is injective. Every actual
degree-zero homogeneous polynomial is its own degree-zero component,
which is the constant polynomial with its original constant coefficient.
Thus every degree-zero element has a genuine constant preimage, proving
surjectivity and the whole ring isomorphism. Polynomial finite generation
also supplies actual finite generation over the degree-zero component
by restriction of scalars. The genuine Mathlib projective structure
morphism to that component's spectrum is therefore proper. Composing
it with the actual spectrum isomorphism of the preceding constant ring
isomorphism proves properness of the original coefficient structure map.
\end{proof}

\begin{theorem}[A value-one normalized chart contains the actual proper center]
\label{projective:native-normalized-value-one-center}
\lean{EnvelopeProperty.Berkovich.Spectrum.isBounded_mvPolynomial_of_C_X,
EnvelopeProperty.Berkovich.RelativeNormalizedProjectiveSpectrum.isBounded_polynomial,
EnvelopeProperty.Berkovich.RelativeNormalizedProjectiveSpectrum.isBounded_coordinateChart_of_val_X_eq_one,
EnvelopeProperty.Berkovich.RelativeNormalizedProjectiveSpectrum.center_nativeHomeomorph_mem_coordinateOpen}
\leanok
\uses{projective:native-polynomial-structure-proper,projective:general-relative-normalized-charts,scheme:proper-center-bounded-affine-criterion}
Let $v$ be a real seminorm on an arbitrary commutative ring $R$, bounded
by one on every coefficient. Let $x$ be a normalized graded polynomial
seminorm with those precise coefficient values, and suppose
$|X_j|_x=1$. Its genuine native projective point has its actual proper
model center in $D_+(X_j)$. Its entire affine coordinate seminorm and
its original whole polynomial seminorm are bounded by one.
\end{theorem}
\begin{proof}
A polynomial seminorm bounded on every original coefficient and variable
is bounded on every polynomial: use true polynomial induction,
ultrametric addition and multiplicativity. The actual normalized
variables have values at most one and the original prescribed coefficients
satisfy the stated bound, so this proves the whole polynomial bound.
On the selected coordinate chart, each genuine variable has ratio value
$|X_i|_x/|X_j|_x\le1$, and every coefficient retains the original value
$v(r)\le1$. The same true whole-polynomial bound proves boundedness of
the entire coordinate ring.

Every real affine section is represented through the genuine coordinate
section-ring isomorphism. The actual normalized/native comparison
preserves each of those entire section values, so all sections on the
selected actual affine open have value at most one. The independently
constructed proper-center theorem then forces the actual center into
that open. Properness was derived above; the chosen center, section
bounds and coordinate evaluations are not assumed certificates.
\end{proof}

\begin{theorem}[Whole homogeneous ideal norm on a value-one chart]
\label{projective:native-homogeneous-ideal-coordinate-norm}
\lean{EnvelopeProperty.Berkovich.RelativeNormalizedProjectiveSpectrum.coordinateChart_val_dehom_of_isHomogeneous,
EnvelopeProperty.Berkovich.RelativeNormalizedProjectiveSpectrum.coordinateChart_val_dehom_le,
EnvelopeProperty.Berkovich.RelativeNormalizedProjectiveSpectrum.idealNorm_coordinateDehom}
\leanok
\uses{projective:general-normalized-charts,projective:native-normalized-value-one-center,lem:general-ideal-pullback}
Under these bounded-coefficient hypotheses, for every original homogeneous
ideal $J$, its intrinsic whole polynomial ideal norm at $x$ equals the
intrinsic norm of its original dehomogenized coordinate ideal at the
actual selected affine chart point. No finite generators are required. The coefficient bound is on the
chosen model ring; it is not a bound on all elements of a nontrivially
valued field.
\end{theorem}
\begin{proof}
For a homogeneous form $F$ of degree $d$, its original dehomogenized
coordinate value is the true fraction value $|F|_x/|X_j|_x^d=|F|_x$.
This is a whole homogeneous-fraction evaluation theorem, rather than an
inference from norms of coordinate generators. For a general polynomial,
decompose its entire finite homogeneous sum. Ultrametricity bounds its
dehomogenized value by the maximum of those actual component values,
which is its original norm by true gradedness. Consequently every
original element of $J$ has dehomogenized value bounded by the original
ideal norm. Actual bounded-coordinate ideal-map/span calculus extends
this inequality to the whole image ideal, giving the first norm bound.

Conversely every homogeneous component of any original $F\in J$ belongs
to the original ideal by genuine homogeneity, so its dehomogenization
lies in the actual image ideal. The original graded evaluation identity
makes $|F|_x$ the maximum of these component values. Each equals its
actual dehomogenized value and is bounded by the whole coordinate ideal
norm. Taking the supremum over every original $F\in J$ proves the
opposite inequality. These arguments identify entire ideal suprema,
not whole polynomial values inferred from selected homogeneous values.
Zero ideals and bottom logarithms are retained.
\end{proof}

\begin{theorem}[Genuine center-stalk model ideal equals the original homogeneous ideal function]
\label{projective:native-homogeneous-model-ideal-representation}
\lean{EnvelopeProperty.ProjBerk.nativeHomogeneousIdealSheaf,
EnvelopeProperty.Berkovich.RelativeNormalizedProjectiveSpectrum.modelIdealNorm_quotient_eq_idealNorm,
EnvelopeProperty.Berkovich.RelativeNormalizedProjectiveSpectrum.modelIdealLog_quotient_eq_idealLog}
\leanok
\uses{projective:native-normalized-value-one-center,projective:native-homogeneous-ideal-coordinate-norm,graded:native-polynomial-quotient-coordinate-ideal,scheme:native-model-ideal-function}
Let $\mathcal J$ be the actual kernel ideal sheaf of the original
homogeneous projective quotient by $J$. At the genuine native point of
$x$, its intrinsic proper-center model ideal norm equals the whole
polynomial ideal norm of $J$ at $x$. Its intrinsic model logarithm equals
the corresponding whole polynomial ideal logarithm. This uses the
original ideal, including powers and nilpotents, with no finite-generator,
noetherian, field, completeness or infinitude hypothesis.
\end{theorem}
\begin{proof}
Choose a normalized coordinate of value one, whose existence is the
actual finite normalization identity. The preceding whole affine bound
puts the genuine proper center in that actual affine open. The true
center-stalk model ideal theorem computes the intrinsic norm on this
affine neighborhood, using the actual ideal sheaf's affine section ideal.
The proved original-ideal coordinate section comparison identifies that
whole ideal, under the true section-ring isomorphism, with the original
dehomogenization image of $J$. Bounded ideal-map calculus and the entire
section-evaluation identity therefore identify this norm with the actual
coordinate ideal norm. The preceding homogeneous-ideal norm comparison
identifies it with the whole original polynomial ideal norm at $x$.
Taking the genuine extended-real logarithm gives the logarithmic
identity, including zero ideals and negative infinity. No radical
replacement or analytic structure-sheaf identification is involved.
\end{proof}
